% NB14a -- Arquitectura de control del turbohelice: palanca de potencia -> control de combustible (gobierna Ng),
% palanca de condicion -> gobernador de helice (gobierna N2 cambiando el paso). Dos lazos casi desacoplados.
\documentclass[tikz,border=4mm]{standalone}
\usepackage{amsmath,amssymb}
\usetikzlibrary{arrows.meta,calc,positioning}
\begin{document}
% Estilo comun de los diagramas TikZ del modulo VCR (NB10b) -- paleta del proyecto
\definecolor{manivela}{HTML}{2A78D6}
\definecolor{inferior}{HTML}{4A3AA7}
\definecolor{control}{HTML}{1BAF7A}
\definecolor{superior}{HTML}{EB6834}
\definecolor{tinta}{HTML}{3A3A38}
\definecolor{metal}{HTML}{C3C2B7}
\definecolor{perdida}{HTML}{E34948}
\definecolor{fuerza}{HTML}{C0392B}
\tikzset{
  barra/.style={line width=3.2pt, line cap=round},
  pin/.style={circle, fill=white, draw=tinta, line width=0.9pt, inner sep=1.6pt},
  fijo/.style={circle, fill=tinta, inner sep=2.6pt},
  vec/.style={-{Stealth[length=2.6mm]}, line width=0.9pt},
  F/.style={-{Stealth[length=3mm]}, line width=1.2pt, fuerza},
  cota/.style={{Stealth[length=1.8mm]}-{Stealth[length=1.8mm]}, line width=0.5pt, tinta},
  lbl/.style={font=\small, text=tinta},
}
\newcommand{\tierra}[1]{\fill[pattern=north east lines, pattern color=tinta] ($(#1)+(-5,-3)$) rectangle ($(#1)+(5,-6)$);
  \draw[tinta, line width=0.8pt] ($(#1)+(-5,-3)$) -- ($(#1)+(5,-3)$);}

\definecolor{gg}{HTML}{4A3AA7}
\definecolor{caliente}{HTML}{E34948}
\definecolor{pt}{HTML}{EB6834}
\definecolor{helice}{HTML}{1BAF7A}
\begin{tikzpicture}[x=1cm,y=1cm,
  b/.style={draw=tinta, line width=0.9pt, rounded corners=2pt, align=center, font=\footnotesize, minimum height=12mm, text width=30mm, fill=white},
  s/.style={-{Stealth[length=2.2mm]}, line width=1pt},
  tag/.style={font=\scriptsize\itshape, text=tinta!75}]
  \node[b, fill=metal!30, text width=24mm] (pla) at (0,3) {\textbf{Palanca de potencia}\\$\Rightarrow N_{g,ref}$};
  \node[draw, circle, inner sep=1.5pt] (s1) at (2.6,3) {$-$};
  \node[b, fill=caliente!8] (fcu) at (5.6,3) {\textbf{Control de combustible}\\PI sobre $N_g$};
  \node[b, fill=caliente!15, text width=34mm] (lim) at (10.0,3) {l\'imites: $R_{dec}\le w_f/p_3\le R_{acc}$\\ITT $\le 800$ $^\circ$C};
  \node[b, fill=gg!15] (gg) at (14.4,3) {\textbf{Generador de gas}\\c\'amara + TC + compresor};
  \node[b, fill=metal!30, text width=24mm] (cl) at (0,-0.3) {\textbf{Palanca de condici\'on}\\$\Rightarrow N_{2,ref}$};
  \node[draw, circle, inner sep=1.5pt] (s2) at (2.6,-0.3) {$-$};
  \node[b, fill=helice!12] (gov) at (5.6,-0.3) {\textbf{Gobernador}\\(contrapesos + v\'alvula piloto)};
  \node[b, fill=helice!20, text width=34mm] (pas) at (10.0,-0.3) {paso $\beta$ entre tope bajo\\y bandera (84.5$^\circ$)};
  \node[b, fill=pt!15] (pt) at (14.4,-0.3) {\textbf{Turbina libre + h\'elice}\\$I\dot\omega_p=P_{TP}-Q\Omega$};
  \draw[s] (pla) -- (s1); \draw[s] (s1) -- node[above, tag] {$e_{N_g}$} (fcu); \draw[s] (fcu) -- node[above, tag] {$w_f^{*}$} (lim);
  \draw[s] (lim) -- node[above, tag] {$w_f$} (gg);
  \draw[s] (gg.north) -- ++(0,0.6) -| node[pos=0.25, above, tag] {$N_g,\;p_3,\;$ITT} (s1.north);
  \draw[s] (cl) -- (s2); \draw[s] (s2) -- node[above, tag] {$e_{N_2}$} (gov); \draw[s] (gov) -- node[above, tag] {$\dot\beta$} (pas);
  \draw[s] (pas) -- node[above, tag] {$Q(\beta)$} (pt);
  \draw[s] (pt.south) -- ++(0,-0.6) -| node[pos=0.25, below, tag] {$N_2$} (s2.south);
  \draw[s, pt, line width=1.4pt] (gg.south) -- node[right, tag] {gas: $p_{45},\,T_{45}$ $\Rightarrow P_{TP}$} (pt.north);
  \node[tag, text width=15cm, anchor=north west] at (-1.2,-2.1) {Acoplamiento: el gas del generador alimenta la turbina libre, pero el \'unico v\'inculo mec\'anico es aerodin\'amico. Subir la palanca de potencia aumenta $N_g$ y $P_{TP}$; el gobernador responde aumentando el paso para que $N_2$ no cambie, y el par $Q$ es lo que el piloto lee en cabina.};
\end{tikzpicture}
\end{document}
